% !PRINTED 445
\weberpage{445}
4. If $(k), (k'), \ldots, (k^{(h-1)})$ denote all the roots of the
equation
\[
H_m(u) = 0,
\]
then for a variable $t$
\[
H_m(t) = [t - (k)] [t - (k')] \ldots [t - (k^{(h-1)})],
\]
and this function has integral rational coefficients. Hence by (5)
we have:
\[
[H_m(t)]^p \equiv H_m(t^p) \quad (\mathrm{mod}\ p),
\]
and so, if we set $t = (k)$:
\[
[(k)^p - (k)] [(k)^p - (k')] \ldots [(k)^p - (k^{(h-1)})] \equiv 0 \quad (\mathrm{mod}\ p).
\]
If, then, $\mathfrak P$ denotes a prime ideal in $\mathfrak K$ that divides $p$,
at least one of the factors of the left side must be divisible by
$\mathfrak P$, and from this there follows a congruence
\weq{12}{(k)^p \equiv (k') \quad (\mathrm{mod}\ \mathfrak P),}
in which $(k')$ is some one of the class invariants, and it may also
be identical with $(k)$.

5. If first
\[
(D, p) = -1,
\]
then $\mathfrak P$ can in no case be of the first degree, since already
\weq{13}{\begin{aligned}
\sqrt{D}^p &\equiv -\sqrt{D}\\
\sqrt{D}^{p^2} &\equiv \sqrt{D}
\end{aligned} \quad (\mathrm{mod}\ \mathfrak P)}
holds. But if in (12) (according to § 121)
% [sic] The German print mixes parenthesized class invariants with bare k in the next three displays.
\[
(k') = f(k, \sqrt{D}), \quad k = f(k', -\sqrt{D}),
\]
then it follows from (13) that
\[
(k')^p \equiv f(k^p, -\sqrt{D}) \equiv f(k', -\sqrt{D}) \equiv k,
\]
\[
(k)^{p^2} \equiv k \quad (\mathrm{mod}\ \mathfrak P),
\]
and consequently, for all integers $\xi$ of the field $\mathfrak K$,
\weq{14}{\xi^{p^2} \equiv \xi \quad (\mathrm{mod}\ \mathfrak P),}
and $\mathfrak P$ is of the \gesperrt{second degree}.

6. If secondly
\[
(D, p) = +1,
\]
then there are, as was proved in § 118, two opposite classes
$l, l^{-1}$ (which, if $l$ is two-sided, are identical with one another),
which can be represented by forms with the first coefficient $p$;
and if $(k)$ is any class invariant, then by § 120 the invariant
equation for the degree of transformation $p$ is:
\weq{15}{F_p(u, v) = 0}
%CONT>>
